\maketitle
\tableofcontents

% ============================================================
\chapter*{Doc.~308 -- The Cosmic Sector of FFGFT}
\addcontentsline{toc}{chapter}{Doc.~308 -- The Cosmic Sector of FFGFT}
% ============================================================

% ============================================================
\section*{What this is about}
% ============================================================

FFGFT advances an unusual claim: the universe is not expanding.
Cosmic redshift is not a Doppler effect but a geometric property of the
$\xi$-field. Dark energy does not exist as a physical entity. And dark
matter requires no particles — its effects follow from the same geometry
that carries light deflection.

This document assembles the cosmic sector of FFGFT: what has been derived,
what has been observationally confirmed, and where the limits lie.
It addresses readers with a physics background who want to know \emph{what}
the theory says and \emph{why}.

The three sources of degeneracy are named, the two cores of the geometry
are developed, and the status of each step — derivation, structural
argument, posit — is declared.

% ============================================================
\section*{Conventions}
% ============================================================

\noindent\textbf{Legend.}

\medskip
\begin{tabular}{@{}lp{11cm}}
\textbf{[K]}      & \emph{Core derivation} — derived mathematically or
                    numerically from the foundational posits; reproducible. \\[4pt]
\textbf{[S]}      & \emph{Structural claim} — justified but not yet fully
                    proved. \\[4pt]
\textbf{[POSIT]}  & \emph{Posit} — declared assumption, not further derived. \\
\end{tabular}

\medskip
\noindent\textbf{Central quantities.}

\medskip
\begin{tabular}{@{}lp{10.5cm}}
$\xi = 4/30000$          & Only fundamental parameter of FFGFT
                           [POSIT, P33]. \\[4pt]
$\tilde T \cdot m = 1$   & Foundational relation: intrinsic time and
                           mass are inversely coupled. \\[4pt]
$R_H = c/H_0$            & Geometric maximal scale — not an expansion
                           horizon. \\[4pt]
$a_0 = c^2\xi^{10}/(4\bar\lambda_e)$
                         & Transition scale of the inertia regime;
                           below $a_0$ inertia is non-linear
                           ($\tilde T\cdot m=1$ + Unruh);
                           $= 1.033\times10^{-10}$~m/s$^2$. \\[4pt]
$\bar\lambda_e = \hbar/(m_e c)$
                         & Reduced Compton wavelength of the electron. \\
\end{tabular}

% ============================================================
\chapter{The Static Universe and Redshift}
% ============================================================

\section{The fundamental choice}

$\Lambda$CDM describes an expanding universe: galaxies move apart, light
is redshifted by the expansion of space, and the observation of distant
supernovae reveals an accelerated expansion accounted for by a cosmological
constant $\Lambda$.

FFGFT makes a different fundamental choice: the universe is static.
Cosmic redshift arises not from motion but from geometric path elongation
in the $\xi$-field along the light path.

\section{The redshift mechanism [K]}

The FFGFT redshift formula (Doc.~280):
\[
  1 + z = e^{\xi\, x}
  \qquad\Longrightarrow\qquad
  z \approx \xi\, x \quad\text{(for small }z\text{),}
\]
where $x$ is the path length. Each photon loses energy along its path
through fractal path elongation — proportional to distance, without
wavelength dependence.

\textbf{Achromatic [K]:} The path elongation is purely geometric and
stretches \emph{all} wavelengths by the same factor $e^{\xi x}$.
Finite-element simulations of the geometry confirm this (K6, Doc.~190).
A chromatic redshift would destroy the coherence of the CMB spectral
shape — this is not observed.

\section{The degeneracy of the data [K]}

For the three classical pillars of cosmology — Type Ia supernovae, baryon
acoustic oscillations (BAO), CMB temperature — the following holds:
they cannot in principle distinguish between expansion and static geometric
path stretching, as long as both models produce the same distance-redshift
relation $z(d)$ (Doc.~267).

This is not a weakness of FFGFT. It is a methodological statement:
\emph{these data alone do not decide between the models.}
The decision shifts to theoretical consistency, parameter count and
derivability.

\section{What follows from this}

\textbf{For dark energy:} $\Lambda$ is needed in $\Lambda$CDM to account
for the apparent accelerated expansion. In the static reading there is no
expansion to be accelerated. $\Lambda$ loses its referent in the static
reading — it is not falsified, it becomes objectless.

\textbf{For the Hubble tension:} Different measurement methods for $H_0$
give different values. In the FFGFT reading, $H_0$ does not measure an
expansion rate but a parameter of the $\xi$-geometry. The tension between
early- and late-universe determinations is then a symptom of the wrong
reading, not a fundamental problem.

\textbf{For the $z$-axis:} Every presentation of cosmological data ``over
time'' is formulated in $\Lambda$CDM language. An FFGFT interpretation must
treat $z$ as a geometric distance parameter, not a time axis.

% ============================================================
\chapter{Geometric Gravity: the Two-Halves Synthesis}
% ============================================================

\section{Einstein 1911 and 1915}

Einstein calculated the deflection of light at the Sun in 1911 from time
dilation alone:
\[
  \alpha_{1911} = \frac{2GM}{c^2 b} = 0.875'' .
\]
That was half the measured value. The 1915 theory of GR yielded the full
value $\alpha_{1915} = 1.750''$, because space \emph{and} time curvature
each contribute one half.

FFGFT \emph{explains} this factor of 2 from first structure.

\section{K3 area invariance [K]}

$\tilde T \cdot m = 1$ (time curvature) and fractal path elongation (space
curvature, Doc.~182) are in FFGFT not alternative mechanisms but
\textbf{complementary halves}. The K3 rule — area invariance of the
time-space 2-cell — anchors this:
\[
  L_t \cdot l = \text{const.}
\]
This is formally the reciprocity form of $\tilde T \cdot m = 1$, applied
to the time-space pair.

\textbf{FE verification [K]:} Ray tracing through three configurations
(30001 log nodes, RK4,
\texttt{ffgft\_308\_p44\_stufeB1\_faktor2\_fe.py}):

\begin{center}
\begin{tabular}{lcc}
\toprule
Configuration & $\alpha$ at solar limb & Ratio to $1.75''$\\
\midrule
Time rate only ($\tilde T \cdot m = 1$) & $0.875''$ & $0.500$\\
Path elongation only (space)            & $0.875''$ & $0.500$\\
Both together (K3)                      & $1.750''$ & $1.000$\\
\bottomrule
\end{tabular}
\end{center}

\textbf{PPN test [K]:} K3 is the sole survivor of the T$^4$ selection
computation (four candidates tested,
\texttt{ffgft\_308\_p44\_stufeB2\_kappa\_herleitung.py}):
$\gamma_\text{PPN} = 1$ exactly. Cassini bound
$|\gamma - 1| < 2.3\times10^{-5}$ passed with unlimited margin.

\section{The same structure in the inertia regime: $a_0$ [K]/[S]}

The K3 two-halves structure appears a second time, at cosmological scale.
From Doc.~279:
\[
  \frac{H_0}{c} = \frac{\pi}{2}\,\frac{\xi^{10}}{\bar\lambda_e}
  \qquad\Longrightarrow\qquad
  R_H = \frac{2\,\bar\lambda_e}{\pi\,\xi^{10}}.
\]

$R_H$ is the geometric maximal scale of the compact T$^4$ geometry
(Doc.~182) — not an expansion horizon, valid in the static reading.

\textbf{The cycle doubling [S]:} The K3 2-cell has two cycles. A closed
path that traverses each once measures:
\[
  L = \underbrace{2\pi R_H}_{\text{time cycle}}
    + \underbrace{2\pi R_H}_{\text{space cycle}}
    = 4\pi R_H .
\]

\textbf{Rindler saturation [S]:} The Rindler horizon of an observer with
acceleration $a$ lies at $c^2/a$. When it reaches $L$, the acceleration is
no longer relationally definable. This sets:
\[
  a_\text{bg} = \frac{c^2}{4\pi R_H}
              = \frac{c^2\,\xi^{10}}{8\,\bar\lambda_e}.
\]

\textbf{Inertia from $\tilde T \cdot m = 1$ [S]:} The foundational relation
is a statement about mass. If inertia follows the time field (Unruh
coupling), the force law in closed form is:
\[
  \boxed{a = \sqrt{a_N^2 + a_N\,a_0}}
  \qquad\text{with}\quad
  a_0 = 2\,a_\text{bg}
      = \frac{c^2\,\xi^{10}}{4\,\bar\lambda_e}
      = 1.033\times10^{-10}~\text{m/s}^2.
\]

The $2\pi$ of the geometry cancels exactly against the $\pi/2$ from
Doc.~279. The rational factor~4 remains:
\[
  4 = \underbrace{2}_{\text{K3: time and space cycle}}
    \times
    \underbrace{2}_{\text{Unruh: }a_0 = 2a_\text{bg}}.
\]
No step contains a free parameter.

% ============================================================
\chapter{The Dark Sector}
% ============================================================

\section{Dark energy: not a physical object}

$\Lambda$ arises in $\Lambda$CDM as an answer to the apparent accelerated
expansion. In the static reading this expansion does not exist. $\Lambda$
thereby loses its referent — it is not falsified by the reading change,
it becomes objectless.

Three conditions speak against it:

\begin{enumerate}
  \item \textbf{Only from degenerate data.} SNe~Ia, BAO, CMB — exactly
        the datasets between which both readings are indistinguishable.
  \item \textbf{No independent anchor.} The only theoretical prediction
        (QFT vacuum energy) misses the value by $10^{122}$.
  \item \textbf{Objectless under reading change.} Not refuted, but
        without referent.
\end{enumerate}

That is the status: reading artefact. Not ``dark energy does not exist''
— that would be a statement within the expansion reading — but: it is a
reading-internal bookkeeping quantity.

\section{Dark matter: a regime, not a substance}

The findings — rotation curves, gravitational lensing, Bullet Cluster —
are reading-independent. They do not disappear under a reading change.
But they require no new ingredient: they follow from the same force law
that $\tilde T \cdot m = 1$ together with K3 geometry delivers.

Dark matter in FFGFT is not a substance and not an open problem.
It is a name for a regime.

\textbf{The force law in the galactic regime [K]/[S]:}
\[
  a = \sqrt{a_N^2 + a_N\,a_0} .
\]
For $a_N \gg a_0$ (inner regions): $a \to a_N$ (Newton).
For $a_N \ll a_0$ (outer regions): $a \to \sqrt{a_N a_0}$,
hence $v^4 = G M_b a_0$ — the baryonic Tully-Fisher relation.

\textbf{In FFGFT, dark matter does not denote a missing mass component
but the regime below $a_0$ in which inertia is non-linear. The missing
mass is not missing — inertia is different there.}

\textbf{Results [K]
(\texttt{ffgft\_308\_p44\_stufeA\_a0\_kette.py}):}

\begin{center}
\begin{tabular}{lcccc}
\toprule
Galaxy & Type & $v_\text{Newton}$ & $v_\text{FFGFT}$ & $v_\text{obs}$ \\
\midrule
DDO~154    & dwarf, gas-dominated & 18.5 & 47.1 & 47.0~km/s \\
NGC~3198   & spiral               & 50.4 & 121  & 150~km/s \\
Milky Way  & spiral               & 122  & 180  & 220~km/s \\
\bottomrule
\end{tabular}
\end{center}

DDO~154 is the cleanest test (gas-dominated, $M_b$ best known):
$v_\text{FFGFT}/v_\text{obs} = 1.001$. Deviations for NGC~3198 and the
Milky Way ($\sim$18\%) lie within the $M_b$ uncertainty ($\sim$30\%,
corresponding to $\sim$8\% in $v$).

\textbf{Bullet Cluster [K]
(\texttt{ffgft\_308\_p44\_stufeC\_bullet\_cluster.py}):}
Two-clump FE projection (Clowe+2006 parameters): lensing peak 10~kpc from
the galaxies, 240~kpc from the gas. Passed.
Cause: $M_\text{gal}/M_\text{gas} = 10$ dominates at instantaneous coupling.

\section{Asymmetry in comparison}

$\xi$ is cross-anchored — it carries the lepton mass ladder, the Koide
scalar, the $\alpha$ reference, $K_\text{frak}$ and $k^*$ (Doc.~275).
It works in five foreign sectors before the cosmological one.

$\Lambda$ works nowhere except its own fit.

That is the difference: not ``true vs.\ false'' but
``cross-anchored vs.\ isolated''.

% ============================================================
\chapter{Summary: Picture of the Cosmic Sector}
% ============================================================

\begin{center}
\begin{tabular}{lll}
\toprule
Quantity & FFGFT & Status \\
\midrule
Cosmic redshift & $1+z = e^{\xi x}$ & [K] \\
Achromatic & geometric, FE confirmed & [K] \\
Degeneracy SNe/BAO/CMB & between both readings & [K] \\
$\Lambda$ (DE) & reading artefact, no referent & [K] \\
Hubble $H_0$ & $c\pi\xi^{10}/(2\bar\lambda_e)$ & [K] \\
Light deflection & $0.875'' + 0.875'' = 1.750''$ & [K] \\
$\gamma_\text{PPN} = 1$ & K3 sole survivor & [K] \\
$a_0 = c^2\xi^{10}/(4\bar\lambda_e)$ & closed chain & [K] \\
Rotation curves & DDO~154 ratio 1.001 & [K] \\
Bullet Cluster & peak 10~kpc from galaxies & [K] \\
K3 cycles $4\pi R_H$ & time $+$ space cycle equal & [S] \\
Rindler saturation & horizon meets maximal scale & [S] \\
$\xi = 4/30000$ & magnitude & [POSIT] \\
K3 forward proof & 2-cell vs.\ 4-volume & open \\
14\% to $a_0^\text{emp}$ & literature range $1.1$--$1.3\times10^{-10}$ & open \\
\bottomrule
\end{tabular}
\end{center}

\medskip
The cosmic sector of FFGFT rests on a single mechanism: $\tilde T \cdot
m = 1$ together with K3 area invariance produces light deflection,
geometric gravity and achromatic redshift. $a_0$ follows from the same
chain without a free parameter.

Dark energy is not a physical object — it is a construct of the expansion
reading. Dark matter as particles is not required — the effects follow from
the geometry.

% ============================================================
\chapter{References}
% ============================================================

\textbf{Internal documents:}
Doc.~025 (cosmology overview);
Doc.~026 (FE geometry);
Doc.~064 (redshift foundations);
Doc.~182 (maximal scale);
Doc.~267 (cosmological degeneracy);
Doc.~279 ($H_0$ from $\xi^{10}$ chain, K6);
Doc.~280 (achromatic redshift);
Doc.~275 ($k^*$ recursion);
Doc.~190 (P33, P44, K6).

\medskip
\noindent\textbf{Verification scripts (all inputs declared, no fits):}
\texttt{ffgft\_308\_p44\_stufeA\_a0\_kette.py}
(chain $\xi \to a_0$, rotation curves);
\texttt{ffgft\_308\_p44\_stufeB1\_faktor2\_fe.py}
(factor-2 gate, 30001 log nodes, RK4);
\texttt{ffgft\_308\_p44\_stufeB2\_kappa\_herleitung.py}
(selection computation K3);
\texttt{ffgft\_308\_p44\_stufeC\_bullet\_cluster.py}
(Bullet Cluster projection).

\medskip
\noindent\textbf{External references:}
McGaugh, Lelli, Schombert 2016, PRL 117, 201101 (RAR/$a_0$);
Clowe et al.\ 2006, ApJ 648, L109 (Bullet Cluster).
